% ----------------------------------------------------------
% Introdução
% ----------------------------------------------------------
\chapter{Introdução}

A automação industrial e o conceito de Indústria 4.0 trouxeram uma demanda crescente por sistemas autônomos, destacando-se os robôs móveis autônomos (AMR). Essas máquinas são fundamentais para otimizar processos logísticos e industriais, operando com segurança e flexibilidade sem a necessidade de intervenção humana direta \cite{Anand_2024_intro, IFR2025, braun2022robot}. A transição para fábricas inteligentes exige não apenas robôs capazes de executar tarefas repetitivas, mas sistemas que possuam percepção sensorial apurada e capacidade de tomada de decisão local. 

Contudo, o desenvolvimento de sistemas robóticos móveis enfrenta desafios relacionados à complexidade da integração entre hardware e software. No ambiente acadêmico, muitas soluções comerciais operam como "caixas-pretas", com protocolos proprietários e hardware fechado. Isso limita o aprendizado prático de controle de baixo nível, cinemática e percepção sensorial \cite{patel2023open, jolles2022open}. Diante dessa barreira, o uso de plataformas de hardware aberto e sistemas operacionais de código fonte livre torna-se essencial para a formação de novos profissionais e pesquisadores, permitindo a exploração profunda de cada camada do sistema.

O \textit{Robot Operating System} (ROS) surge como um padrão de fato para mitigar esses desafios, oferecendo um framework aberto, modular e colaborativo para o desenvolvimento robótico \cite{macenski2022robot}. A arquitetura do ROS 2, em particular, promove a modularidade através de nós independentes que se comunicam via tópicos, serviços e ações. Isso permite que pesquisadores reutilizem algoritmos consolidados de localização e controle, focando na inovação de suas aplicações específicas sem a necessidade de reinventar tecnologias existentes.

Nesse contexto, este artigo descreve o desenvolvimento de uma plataforma robótica seguidora de linha, adaptada a partir de um robô aspirador comercial (iRobot Roomba). O projeto visa preencher a lacuna entre os modelos matemáticos teóricos e a execução física real \cite{vega2025bridging}, utilizando controle Proporcional-Integral-Derivativo (PID) para a manutenção da trajetória. A arquitetura proposta emprega um Raspberry Pi 4 Model B, rodando Ubuntu 22.04 e o \textit{middleware} ROS 2 Humble \cite{kortelainen2023short}. 

A principal contribuição deste trabalho reside na integração entre o ambiente de simulação e o protótipo físico, utilizando a abordagem de \textit{Edge Computing} (Computação de Borda). Ao centralizar a leitura direta de sensores infravermelhos e o processamento de controle na placa central, elimina-se a latência de microcontroladores intermediários. Assim, a plataforma fornece uma solução acessível, robusta e escalável, preparando possibilidades de aplicações futuras, como em logística hospitalar, almoxarifados inteligentes e até mesmo em outras áreas.